
En
effet, si @@M0@@, @@M1@@ est évidemment un @@M2@@-faisceau. Si @@M3@@, par 2.4, @@M4@@ puisque @@M5@@. Soit @@M6@@ une section hyperplane de @@M7@@ et soit @@M8@@ une suite exacte définie par @@M9@@. Si @@M10@@, par 1.7, @@M11@@; donc @@M12@@ est un @@M13@@-faisceau. Si @@M14@@, on peut prendre @@M15@@ intègre grâce au théorème de Bertini et on peut ainsi supposer par récurrence que @@M16@@ est un @@M17@@-faisceau; d'où l'assertion.

\textbf{Lemme 2.7.} \emph{Si @@M18@@ et @@M19@@, alors} :
\begin{enumerate}[label=(\roman*)]
\item \emph{Pour tout @@M20@@, @@M21@@}
\item \emph{@@M22@@.}
\end{enumerate}

Soit @@M23@@ une section hyperplane intègre de @@M24@@. Par 1.7 et 2.5, @@M25@@ est un @@M26@@-faisceau à polynôme de Hilbert @@M27@@. Donc par 1.11, @@M28@@ et @@M29@@ est @@M30@@-régulier.

Considérons la suite exacte suivante :
@@M31@@
Pour @@M32@@, on a @@M33@@ car @@M34@@ par 1.3 (i); pour @@M35@@, on a @@M36@@ car @@M37@@ par 2.4. Enfin, pour tout @@M38@@, @@M39@@ car il existe une injection @@M40@@, et @@M41@@ car @@M42@@; donc pour @@M43@@, on a @@M44@@. Comme conséquence, on en conclut que pour tout @@M45@@,
@@M46@@

